\section*{T4.5 — DMP Comparison: Initial vs Final}

\noindent\textbf{Initial DMP (Part 2):} \textit{[https://doi.org/10.70124/nf9qd-q7b67]}\\
\textbf{Final DMP:} \url{https://doi.org/10.70124/21cgp-a9396}

\bigskip

\subsection*{Question 1 — What did the initial DMP promise that turned out to be unrealistic?}

\subsubsection*{DBRepo as a straightforward programmatic data layer}

The initial DMP described the experiment's data infrastructure as a clean API-driven
workflow where the schema would be created and data loaded via the DBRepo REST API
using the Python SDK. In practice, the SDK (\texttt{dbrepo==1.13.3}) was severely
misaligned with the test server. The \texttt{create\_table(dataframe=...)} method
mapped \texttt{int64} columns to \texttt{BOOLEAN} type because it inferred types from
Python values rather than numpy dtypes --- meaning a seed value of \texttt{0} was
interpreted as \texttt{False}. The \texttt{pd.StringDtype()} extension type was not
recognised at all, defaulting silently to \texttt{text}. The \texttt{primary\_key}
parameter did not exist on \texttt{CreateTableColumn}. Foreign key constraints declared
in \texttt{schema.sql} could not be enforced through the SDK. None of these limitations
are discoverable without attempting the actual implementation --- an initial DMP written
before touching the tools could not have anticipated them.

\subsubsection*{Exact reproducibility of the API reimplementation}

The initial DMP would have described the API reimplementation as producing identical
results to the original local-file version, as T2.6 requires. This proved unreachable.
The DBRepo database stored \texttt{value\_start\_mln} as \texttt{DECIMAL(12,3)},
introducing rounding at the third decimal place. Additionally, 17 rows with unknown
\texttt{start\_value\_eur} could not be stored as NULL because the column was created
as NOT NULL during table setup --- they were stored as \texttt{0.0} instead. These 17
rows were handled by \texttt{SimpleImputer(strategy=`median')} in the original (filling
NaN with the median) but treated as genuine zero-value observations in the API version,
shifting model predictions. The final RMSE differed by 7\% (8.0957 vs.\ 8.6637) and
MAE by 4\% (5.5635 vs.\ 5.7851). Results were considered equivalent within an 8\%
tolerance, but claiming identical results in the initial DMP was never achievable given
how database storage interacts with ML preprocessing.

\subsubsection*{Schema creation as a simple scripted step}

The initial DMP described schema creation as a one-time programmatic step. In practice,
tables had to be deleted and recreated multiple times --- using DataFrames with various
dtype strategies, using \texttt{CreateTableColumn} with explicit \texttt{ColumnType}
values, and ultimately using the DBRepo UI for tables the SDK could not create
correctly. The \texttt{season} table, whose rows are pre-inserted by the schema SQL,
was incorrectly included in upload pipelines causing PRIMARY KEY duplicate errors. The
fact that DBRepo's \texttt{get\_tables()} returns lightweight \texttt{TableBrief}
objects with no column metadata --- requiring a separate \texttt{get\_table()} call per
table --- meant that any notebook assuming table metadata was immediately accessible
after listing would fail.

\bigskip

\subsection*{Question 2 — What did writing the final DMP actually require that the initial one glossed over?}

\subsubsection*{View architecture as an engineering decision}

The initial DMP would have stated that views would expose query-ready feature tables for
the ML pipeline. The final version required a specific architectural decision that could
not have been anticipated: DBRepo's SDK mapper internally calls \texttt{numpy.array()}
on each joined table's column list. When joined tables have different numbers of
columns, the array becomes inhomogeneous and raises a \texttt{ValueError}. This meant
JOIN-based views --- the natural implementation --- were impossible through the SDK. The
working solution was to create single-table views without joins and reconstruct the
denormalised dataset locally in pandas after retrieving each view via the REST API. This
decision had downstream consequences: lookup views (\texttt{vw\_player\_lookup},
\texttt{vw\_club\_lookup}, etc.) returned HTTP 500 errors when the underlying tables
were empty, requiring lookup data to be fetched via \texttt{get\_table\_data()} instead
of the view API. None of this appears in any initial DMP --- it represents knowledge
gained exclusively through implementation.

\subsubsection*{Column type precision in database storage}

The final DMP documents specific storage decisions that the initial version could not
have specified. \texttt{DECIMAL} columns require both \texttt{size} and \texttt{d}
parameters in the DBRepo UI --- specifying only one triggers a \texttt{BAD\_REQUEST}
validation error. \texttt{BOOLEAN} columns must be uploaded as integer \texttt{1}/\texttt{0}
not as Python \texttt{True}/\texttt{False} strings. Columns defined as
\texttt{SMALLINT NULL} in the schema become \texttt{float64} in pandas when they
contain NaN values --- which DBRepo then maps to \texttt{DECIMAL} type rather than
integer, causing upload rejections for values such as \texttt{-1}. Each of these
required a specific fix: \texttt{pd.Int64Dtype()} for nullable integers,
\texttt{\{1.0: 1, 0.0: 0\}} mapping for boolean columns, \texttt{lineterminator=`\textbackslash n'}
and plain \texttt{utf-8} encoding (not \texttt{utf-8-sig}) to avoid BOM characters that
corrupted the DBRepo UI parser. The initial DMP described data loading in a single
sentence; the final version reflects weeks of iterative debugging.

\subsubsection*{Provenance documentation scope}

T2.6 requires documenting the API base URL, all endpoints used, and the authentication
method in the README. The initial DMP would have described this as a brief technical
note. The final implementation required documenting: the exact database ID, seven view
names and their purposes, the distinction between fact views (via
\texttt{get\_view\_data()}) and lookup tables (via \texttt{get\_table\_data()} due to a
server bug), the SDK version (\texttt{dbrepo==1.13.3}), the column rename mapping
between DBRepo column names and the original Excel column names expected by the ML
pipeline, the reconstruction of \texttt{start\_value} in EUR from the stored
\texttt{value\_start\_mln} column, and a quantified equivalence check with explanation
of the measured difference. An audit trail CSV (\texttt{transfer\_data\_from\_api.csv})
was also produced to make the exact API input to the model reproducible. The initial DMP
described none of this specificity because it predates any knowledge of what the actual
implementation would require.

\subsubsection*{Three separately justified licence categories}

The initial DMP likely assigned a single licence to the project. The final DMP required
treating three artefact categories independently. Input data is governed by the original
publishers' CC~BY~4.0 licence (Mendeley Data) --- the group is a re-publisher, not a
rights holder, and this distinction had to be reflected explicitly in DBRepo metadata
and the \texttt{source\_dataset} table. Source code required a separate open-source
licence (MIT) with a written justification of its compatibility with CC~BY~4.0 and a
\texttt{LICENSE} file in the repository root. Generated outputs (models, predictions,
metrics) required a third licence choice (CC~BY~4.0) stated in every TUWRD deposit
record. The initial DMP would not have separated these categories or addressed the
compatibility analysis between them.

\subsubsection*{Data quality issues that only emerged during loading}

The initial DMP described the source datasets as openly available and suitable for
direct use. 17 rows with missing \texttt{start\_value\_eur} that could not be stored as
NULL due to a column constraint applied during table creation; 6 missing
\texttt{instagram\_followers\_mln} values filled with \texttt{0.0}; and columns such as
\texttt{club\_performance}, \texttt{relegation}, and \texttt{success\_or\_not} that are
genuinely NULL for four of the five seasons because they were only collected in 2019.
Each required a documented decision --- the initial DMP treated data loading as a
straightforward ingestion step.
